\begin{table}[h]
\centering
\caption{Algorithmic Redistricting vs Enacted 2020 Congressional Plans}
\label{tab:enacted_comparison}
\begin{tabular}{lcccccc}
\hline
State & \multicolumn{2}{c}{Enacted 2021-22} & \multicolumn{2}{c}{Baseline METIS} & \multicolumn{2}{c}{Edge-Weighted} \\
 & MM & PP & MM & PP & MM & PP \\
\hline
AL & 1 & 0.185 & 0 & 0.234 & 1 & 0.147 \\
GA & 4 & 0.210 & 5 & 0.204 & 5 & 0.204 \\
LA & 1 & 0.220 & 1 & 0.211 & 1 & 0.211 \\
MS & 2 & 0.240 & 2 & 0.284 & 2 & 0.284 \\
SC & 1 & 0.200 & 0 & 0.238 & 1 & 0.117 \\
\hline
\hline
\textbf{Mean} & & 0.211 & & 0.234 & & 0.193 \\
\textbf{Improvement} & & -- & & +11.1\% & & +-8.6\% \\
\hline
\end{tabular}
\vspace{0.2cm}
\begin{minipage}{0.9\textwidth}
\footnotesize
\textbf{Notes:} 
Enacted plans are 2021-2022 congressional districts adopted by state legislatures. 
MM = Majority-minority districts (50\%+ Black population). 
PP = Average Polsby-Popper score (higher = more compact). 
Baseline METIS uses uniform edge weights (demographic-blind). 
Edge-weighted METIS uses demographic edge weighting for VRA optimization. 
Improvement percentages show compactness gain vs enacted plans. 
Enacted compactness scores from Princeton Gerrymandering Project and court documents.
\end{minipage}
\end{table}